\chapter{Implementación y trazabilidad}

\section{Entorno de software}
El proyecto se implementó en Python, con PyTorch para diferenciación automática y redes neuronales, SciPy para integración ODE, Nilearn y NiBabel para fMRI, NumPy/Pandas para manipulación de datos, scikit-learn para métricas y Matplotlib para visualización \citep{Paszke2019}. Los entrenamientos PINN se ejecutaron con precisión de 64 bits y GPU cuando estuvo disponible.

\section{Pipeline de trece notebooks}
El desarrollo se organizó como una progresión metodológica, no como un único notebook final:
\begin{enumerate}
  \item verificación de datos HCP y extracción de señales;
  \item GLM basal real;
  \item generación y validación del conjunto sintético;
  \item GLM sintético;
  \item MLP sintética;
  \item piloto PINN global;
  \item reformulación PINN inversa de ventana;
  \item lote de 60 PINN sintéticas;
  \item comparación estadística sintética;
  \item preparación y control de calidad real;
  \item piloto de modelos reales;
  \item lote de ocho evaluaciones reales;
  \item análisis final y exportación de tablas y figuras.
\end{enumerate}

La secuencia permite reconstruir por qué se modificó el método. Conservar el piloto fallido es relevante porque evita presentar la arquitectura final como una decisión arbitraria tomada después de ver resultados.

\section{Código modular}
Además de los notebooks, el repositorio contiene módulos reutilizables para:
\begin{itemize}
  \item ecuaciones Balloon--Windkessel y simulación;
  \item conversión de eventos a estímulo;
  \item cálculo BOLD desde estados;
  \item métricas de regresión;
  \item corrección de Holm y tamaño de efecto pareado;
  \item control de semillas.
\end{itemize}

Se incluyeron pruebas fisiológicas básicas: equilibrio sin estímulo y respuesta positiva ante un evento. Las dos pruebas pasaron en Python 3.11. Estas verificaciones no sustituyen una batería completa de validación, pero detectan errores fundamentales de integración y estado basal.

\section{Configuración y reproducibilidad}
Los hiperparámetros se concentraron en archivos JSON y constantes documentadas. Los lotes guardaron resultados después de cada experimento y podían reanudarse desde un archivo de estado. Esta decisión fue necesaria porque 60 PINN sintéticas demandaron aproximadamente dos minutos cada una, alrededor de dos horas de cómputo total.

\begin{lstlisting}[language=bash,caption={Comandos mínimos para instalar y verificar el repositorio.}]
python -m venv .venv
# activar el entorno
pip install -r requirements.txt
python -m pytest
\end{lstlisting}

\begin{lstlisting}[language=bash,caption={Orden conceptual de ejecución.}]
# 1. preparar datos reales y sinteticos
# 2. ejecutar GLM y MLP
# 3. entrenar PINN piloto y PINN corregida
# 4. ejecutar lotes sintetico y real
# 5. consolidar resultados
jupyter notebook notebooks/12_analisis_final_experimentos.ipynb
\end{lstlisting}

\section{Datos no distribuidos}
Los volúmenes fMRI HCP crudos no se incluyen en GitHub por tamaño y condiciones de acceso. El repositorio documenta su estructura esperada y contiene resultados derivados livianos. Tampoco se versionan modelos binarios grandes ni historiales redundantes.

\section{Control de versiones}
El repositorio público conserva código, notebooks, configuraciones, figuras, tablas, informe y pruebas. Cada actualización del informe y la presentación se registra mediante commits. Esto aporta trazabilidad entre el código que generó resultados y las afirmaciones incluidas en la entrega.
